% ==================================================================
% 5. Limitations. Doi khung 2026-07-26: muc "augmentation control is missing"
% DA XOA vi da chay doi chung do (no chinh la ket qua chinh bay gio).
% file: sections/14_limitations.tex
% ==================================================================
\section{Limitations}
\label{sec:limitations}

\subsubsection{The corruption model is synthetic, and it is the load-bearing assumption.}
VisDrone-Shake is built from the known mechanisms of rotorcraft image degradation,
but its parameters bracket flight rather than being calibrated against accelerometer
spectra measured on our own airframe. Every result in this paper is therefore a
statement about a \emph{model} of vibration, and we make no claim about the
magnitude of the effect on any particular aircraft. Only measurement can settle
that: flying the same detector with a synchronised IMU log, fitting $\mathcal{T}$
to the accelerometer spectrum, and re-evaluating. That is the single most valuable
extension of this work, and until it is done the ordering of the models under
vibration is better trusted than the absolute numbers.

\medskip\noindent\textbf{One dataset, one corruption family.}
Every number here comes from the person subset of a single benchmark. VisDrone is
the standard choice for this task and its instance statistics are the ones our
detector targets, but a corruption robustness result measured on one dataset is a
statement about that dataset until it is repeated. We also model one corruption
family, the one rotors produce; nothing here speaks to rain, fog, low light or
motion of the subject rather than the camera, all of which a fielded UAV also meets.

\medskip\noindent\textbf{The scope of the regime is bounded on one side only.}
The gain reproduces on both of our detectors and is absent on the one outside
architecture we tried. Section~\ref{sec:generality} rejects pretraining, object
scale and remaining headroom as explanations, which narrows the cause to the
architecture or the loss without identifying it, and one non-responding model
cannot. We have therefore demonstrated a large effect within one detector family
and one negative instance outside it, with the boundary unknown. A practitioner
should treat the regime as worth trying rather than as established for their
detector, and should run the two controls of Table~\ref{tab:main} when they do.

\medskip\noindent\textbf{Cost, tuning and statistical power.}
Why the consistency term underperforms is measured but not explained: we offer the
reading that the detection loss never reaches the corrupted image, without testing
it, so take the ordering as measured and the reason as conjecture. Clean accuracy
falls by $9.7\%$, roughly a third of that from the finetuning itself, and we did not
try to recover the rest by annealing the corruption probability late in training.
We report no sweep over $p$ or the severity distribution, so those values are
reasonable rather than tuned. Each configuration is trained once, so our error bars
capture evaluation but not training variance. Frames are processed independently,
so the temporal redundancy vibration itself creates, successive frames blurring
along different directions, is available only to the downstream tracker.
